% =============================================================================
% Annexes
% =============================================================================
\appendix
\chapter{Annexe A — Liste complète des features}
\label{ann:features}

\begin{longtable}{|p{4.5cm}|p{2cm}|p{7cm}|}
  \caption{Variables du modèle classique (feature\_order)} \\
  \hline
  \textbf{Feature} & \textbf{Type} & \textbf{Description} \\
  \hline
  \endfirsthead
  \hline
  \textbf{Feature} & \textbf{Type} & \textbf{Description} \\
  \hline
  \endhead
  montant & float & Montant du crédit (TND) \\
  \hline
  duree\_mois & int & Durée en mois \\
  \hline
  dti & float & Debt-to-income ratio \\
  \hline
  revenu\_mensuel & float & Revenu mensuel client \\
  \hline
  age & int & Age du client \\
  \hline
  kyc\_score & float & Score KYC 0--100 (calculé) \\
  \hline
  profession\_enc & int & Profession encodée 0--7 \\
  \hline
  cycle\_enc & int & Cycle crédit encodé \\
  \hline
  objet\_enc & int & Objet crédit encodé \\
  \hline
  avg\_retard & float & Retard moyen (jours) \\
  \hline
  max\_retard & float & Retard maximum \\
  \hline
  std\_retard & float & Ecart-type retards \\
  \hline
  n\_payments & int & Nombre d'échéances \\
  \hline
  n\_late & int & Paiements en retard \\
  \hline
  pct\_late & float & Proportion retards \\
  \hline
  n\_en\_retard & int & Retards sévères ($\geq$90j) \\
  \hline
  n\_transactions & int & Nombre transactions \\
  \hline
  n\_suspect & int & Transactions suspectes \\
  \hline
  avg\_tx\_amount & float & Montant moyen transaction \\
  \hline
  total\_depot & float & Total dépôts \\
  \hline
  total\_retrait & float & Total retraits \\
  \hline
  total\_remboursement & float & Total remboursements \\
  \hline
  total\_transfert & float & Total transferts \\
  \hline
  ratio\_retrait\_depot & float & Ratio retrait/dépôt \\
  \hline
  max\_risk\_relation & float & Risque relation max \\
  \hline
  avg\_risk\_relation & float & Risque relation moyen \\
  \hline
  n\_relations & int & Nombre de relations \\
  \hline
  n\_garant & int & Nombre de garants \\
  \hline
\end{longtable}

\chapter{Annexe B — Guide d'installation}
\label{ann:install}

\section{Prérequis}

\begin{itemize}[leftmargin=*]
  \item Python 3.11+
  \item Node.js 18+
  \item Git
  \item Ollama (optionnel, pour LLM)
  \item 8 Go RAM minimum
\end{itemize}

\section{Installation backend}

\begin{lstlisting}[language=bash,caption={Setup Python}]
cd Talys_pfe
python -m venv venv
venv\Scripts\activate        # Windows
pip install -r requirements.txt
python generate_data.py
python -m src.features.engineering
python -m src.models.train
python -m src.kyc.score
python -m src.models.sequential.train
python -m src.models.graph.train
python -m src.rag.index
uvicorn src.api.main:app --reload --port 8000
\end{lstlisting}

\section{Installation frontend}

\begin{lstlisting}[language=bash,caption={Setup React}]
cd frontend
npm install
npm run dev
\end{lstlisting}

\chapter{Annexe C — Extrait REF-08}
\label{ann:ref08}

\begin{tcolorbox}[colback=codebg,colframe=codeframe,title=REF-08 — Extrait]
\textbf{1.1 Modèles d'Intelligence Artificielle}\\
Objectif: estimer le score KYC ainsi que l'évaluation du profil de risque crédit des clients.

\textbf{1.1.2 Modèles de Scoring Classiques}\\
- Régression Logistique\\
- Forêt Aléatoire\\
- XGBoost

\textbf{1.1.3 Modèles Séquentiels}\\
- LSTM, GRU

\textbf{1.1.4 Modèles de Graphes}\\
- GNN, GraphSAGE

\textbf{1.1.5 Sorties}\\
Score KYC, Score risque crédit, Probabilité de défaut, Explication du score.

\textbf{1.2 Agent Conversationnel}\\
LangChain, tools ML, mémoire conversationnelle.
\end{tcolorbox}

\chapter{Annexe D — Code orchestrateur LangGraph}
\label{ann:langgraph}

\begin{lstlisting}[caption={Extrait CreditAgentOrchestrator — construction graphe},label={lst:orch}]
def _build_graph(self):
    g = StateGraph(AgentState)
    g.add_node("ExtractCIN", self._node_extract_cin)
    g.add_node("DetectIntent", self._node_detect_intent)
    g.add_node("ChoosePath", self._node_choose_path)
    g.add_node("RunClassic", self._node_run_classic)
    g.add_node("RunSequential", self._node_run_sequential)
    g.add_node("RunGraph", self._node_run_graph)
    g.add_node("RunFullReport", self._node_run_full_report)
    g.add_node("RunCompareModels", self._node_run_compare_models)
    g.add_node("RetrieveDocs", self._node_retrieve_docs)
    g.add_node("WriteReport", self._node_write_report)
    g.add_node("Respond", self._node_respond)
    g.add_edge(START, "ExtractCIN")
    g.add_edge("ExtractCIN", "DetectIntent")
    g.add_edge("DetectIntent", "ChoosePath")
    g.add_conditional_edges("ChoosePath", self._route_from_state, {...})
    g.add_edge("RunClassic", "RetrieveDocs")
    # ... autres routes vers RetrieveDocs
    g.add_edge("RetrieveDocs", "WriteReport")
    g.add_edge("WriteReport", "Respond")
    g.add_edge("Respond", END)
    return g.compile()
\end{lstlisting}

\chapter{Annexe E — Règles métier microfinance}
\label{ann:regles}

\begin{tcolorbox}[colback=codebg,colframe=codeframe,breakable]
\textbf{Niveaux de risque}\\
FAIBLE : probabilité faible, KYC conforme, peu de retards.\\
MODERE : vigilance accrue, suivi renforcé.\\
ELEVE : retards répétés, décision humaine obligatoire.

\textbf{Recommandations}\\
FAIBLE $\rightarrow$ validation sous surveillance standard.\\
MODERE $\rightarrow$ conditions renforcées, revue à 3 mois.\\
ELEVE $\rightarrow$ refus ou restructuration comité crédit.

\textbf{Conformité}\\
Données synthétiques en démo. Rapport IA n'altère pas les scores ML.
\end{tcolorbox}

\chapter{Annexe F — Diagrammes PlantUML complets}
\label{ann:plantuml}

\section{Cas d'utilisation}

\begin{lstlisting}[language={},caption={PlantUML Use Case}]
@startuml
left to right direction
actor "Analyste Crédit" as Analyste
actor "Manager Risque" as Manager
actor "Administrateur" as Admin
rectangle "Plateforme Scoring" {
  usecase "Analyser client par CIN" as UC1
  usecase "Chat conversationnel" as UC2
  usecase "Comparer modeles" as UC3
  usecase "Generer rapport RAG" as UC4
  usecase "Telecharger PDF/MD" as UC5
  usecase "Administration API" as UC6
}
Analyste --> UC1
Analyste --> UC2
Analyste --> UC4
Manager --> UC3
Admin --> UC6
@enduml
\end{lstlisting}

\section{Diagramme de classes}

\begin{lstlisting}[language={},caption={PlantUML Class Diagram}]
@startuml
package Frontend { class ReactApp; class ApiClient; }
package Backend { class FastAPIApp; }
package Orchestration {
  class CreditAgentOrchestrator
  class AgentState
}
package Scoring {
  class ClassicScoring
  class SequentialScoring
  class GraphScoring
}
package RAG { class RagIndex; class RagRetriever; }
ReactApp --> FastAPIApp
FastAPIApp --> CreditAgentOrchestrator
CreditAgentOrchestrator --> ClassicScoring
CreditAgentOrchestrator --> RagRetriever
@enduml
\end{lstlisting}

\chapter{Annexe G — Glossaire}
\label{ann:glossaire}

\begin{longtable}{|p{3.5cm}|p{10cm}|}
  \caption{Glossaire} \\
  \hline
  \textbf{Terme} & \textbf{Définition} \\
  \hline
  \endfirsthead
  \hline
  \textbf{Terme} & \textbf{Définition} \\
  \hline
  \endhead
  AUC-ROC & Aire sous la courbe ROC, métrique de discrimination \\
  \hline
  CIN & Carte d'identité nationale (identifiant client) \\
  \hline
  CRISP-DM & Méthodologie data mining en 6 phases \\
  \hline
  DTI & Debt-to-Income, ratio d'endettement \\
  \hline
  GraphSAGE & Graph Sample and Aggregate, GNN inductif \\
  \hline
  GRU & Gated Recurrent Unit, RNN \\
  \hline
  ICM & Interface / système core microfinance (SI métier) \\
  \hline
  KYC & Know Your Customer, connaissance client \\
  \hline
  LangGraph & Framework graphe d'états pour agents LLM \\
  \hline
  LSTM & Long Short-Term Memory, RNN \\
  \hline
  LLM & Large Language Model \\
  \hline
  RAG & Retrieval-Augmented Generation \\
  \hline
  TF-IDF & Term Frequency - Inverse Document Frequency \\
  \hline
\end{longtable}

\chapter{Annexe H — Webographie et bibliographie}
\label{ann:bib}

\begin{thebibliography}{99}

\bibitem{fastapi}
FastAPI Documentation. \url{https://fastapi.tiangolo.com/}

\bibitem{langgraph}
LangGraph Documentation. \url{https://langchain-ai.github.io/langgraph/}

\bibitem{sklearn}
Scikit-learn: Machine Learning in Python. Pedregosa et al., JMLR 2011.

\bibitem{xgboost}
Chen, T., Guestrin, C. XGBoost: A Scalable Tree Boosting System. KDD 2016.

\bibitem{graphsage}
Hamilton, W., et al. Inductive Representation Learning on Large Graphs. NeurIPS 2017.

\bibitem{hochreiter}
Hochreiter, S., Schmidhuber, J. Long Short-Term Memory. Neural Computation, 1997.

\bibitem{cho2014}
Cho, K., et al. Learning Phrase Representations using RNN Encoder-Decoder. EMNLP 2014.

\bibitem{lewis2020}
Lewis, P., et al. Retrieval-Augmented Generation for Knowledge-Intensive NLP. NeurIPS 2020.

\bibitem{crispdm}
Chapman, P., et al. CRISP-DM 1.0 Step-by-step data mining guide. 2000.

\bibitem{faker}
Faker Python Library. \url{https://github.com/joke2k/faker}

\bibitem{ollama}
Ollama — Local LLM Runtime. \url{https://ollama.com/}

\bibitem{react}
React Documentation. \url{https://react.dev/}

\bibitem{pytorch}
PyTorch Documentation. \url{https://pytorch.org/}

\bibitem{ref08}
Talys Consulting. REF-08 — Spécifications Techniques Scoring Crédit. Document interne, 2025.

\bibitem{bale2}
BIS. Basel II — International Convergence of Capital Measurement. BIS, 2004.

\end{thebibliography}
